% ============================================================================
% Section I: Introduction (humanizer-informed revision; original structure retained)
% ============================================================================

\section{Introduction}
\label{sec:introduction}

The growing use of lithium-ion batteries in electric vehicles and energy-storage
systems has increased the demand for efficient fast charging. Fast charging
requires balancing charging speed against thermal stress and degradation. For a
prescribed amount of charge, shortening the charging time requires a higher
average current. However, high charging rates increase polarization and heat
generation and can promote degradation, with the response depending on
temperature and state of charge (SOC)~\cite{refFastChargeReview}.
Charging-protocol design therefore involves selecting how the current varies
across charging stages while reaching a target SOC within prescribed current,
voltage, and temperature limits. The competing objectives of charging time,
thermal response, and battery health motivate a constrained multiobjective
formulation that provides alternative protocols for different operating
preferences~\cite{nref2,refGPA}.

Evaluating these trade-offs is expensive. Physical charging and aging tests
consume cells and laboratory time; in particular, directly measuring cycle
life can require months of testing~\cite{nref3}. Battery models allow candidate
protocols to be assessed in simulation, reducing reliance on physical
experiments during the design stage. Nevertheless, electrochemical models
require repeated numerical solution of coupled transport and reaction
equations, often together with thermal and aging dynamics. Consequently,
replacing experiments with simulation reduces the cost of each evaluation but
does not eliminate the computational burden of searching across many
protocols~\cite{nref7,nref8}. This motivates sample-efficient optimization
methods that require fewer simulator evaluations.

Bayesian optimization (BO) addresses this evaluation cost through sequential
surrogate modeling and acquisition-based sampling. A probabilistic surrogate,
commonly a Gaussian process (GP), estimates both objective values and predictive
uncertainty from evaluated candidates. An acquisition function combines these
estimates to select the next candidate, allowing many alternatives to be
compared using the inexpensive surrogate before one is evaluated by the battery
simulator~\cite{nref6}.

Several studies have established the relevance of BO to fast charging.
Attia \emph{et al.} combined early cycle-life prediction with BO to optimize
fast-charging protocols experimentally, reducing both the duration and number
of tests~\cite{nref3}. Jiang \emph{et al.} investigated BO for minimum-time
charging subject to degradation-related constraints and compared acquisition
functions for multistage protocols~\cite{nref7}. Wang and Jiang extended this
approach to charging-time and cycle-life trade-offs using Tchebycheff
scalarization and constrained BO~\cite{nref8}. More recently,
Zhu \emph{et al.} combined early degradation prediction with multiobjective BO
to consider knee-point capacity and cycle life~\cite{refZhuMOBO}.
GP-assisted evolutionary charging design has also addressed multiple user
preferences while reducing simulation requirements~\cite{refGPA}. Despite
these advances, the search is driven primarily by numerical observations
collected during the optimization process.

Battery-domain knowledge can complement these numerical observations when only
a few evaluations are available. In particular, prior knowledge of how
charging stages affect thermal response and degradation can help identify
useful regions of the protocol space. Prior-guided BO has already been
investigated for battery charging: Jeong \emph{et al.} incorporated a simplified
aging model to guide the search for health-conscious
protocols~\cite{refPriorCharging}. Such model-based prior information, however,
requires an explicit quantitative representation of the available knowledge.
Qualitative knowledge concerning charging-variable meanings, operating
practices, and objective trade-offs is more difficult to incorporate directly
into BO, particularly when candidate selection must account for multiple
competing objectives.

Large language models (LLMs) can interpret such descriptions and generate
candidate suggestions conditioned on the optimization context. Liu \emph{et al.}
introduced LLAMBO and investigated LLM-based warm starting, surrogate modeling,
and candidate sampling, with empirical results on hyperparameter
optimization~\cite{nref12}. In the battery domain, Kuai \emph{et al.} applied
LLM-enhanced BO to electrochemical parameter identification~\cite{refKuaiLLMBO}.
These studies suggest that LLMs can provide task-specific prior information when
numerical observations are scarce. For multiobjective charging design, however,
the guidance must also reflect the trade-off currently being optimized.

These observations suggest two points at which LLM-derived knowledge may be
useful under a limited simulation budget. First, it can improve the initial
design by proposing plausible charging protocols while maintaining sufficient
diversity. Second, it can provide regional guidance during sequential BO
according to the current objective trade-off. Because the scalarization weights
change across iterations, a region that is suitable for one objective trade-off
may be less suitable for another.

This paper proposes \emph{LLM-guided Multiobjective Bayesian Optimization
(LLMBO-MO)} for multistage battery charging. The method combines LLM-informed
initialization with covariance-based regional guidance in a decomposition-based
BO framework. The regional-guidance component is inspired by the region-lifted
preference mechanism of LGBO~\cite{ref5} and is adapted here to the
weight-conditioned multiobjective setting. We refer to this adaptation as
Region-Lift. The main contributions are as follows:

\begin{itemize}

    \item We develop an LLM-informed initialization strategy that uses
    charging-variable definitions, operating constraints, and battery-domain
    knowledge to propose initial protocols. Candidate screening and
    diversity-aware selection are combined with random valid protocols to
    construct the initial simulation design.

    \item We adapt region-lifted preference guidance to decomposition-based
    multiobjective BO by conditioning the regional preference on the active
    scalarization weight and restricting its influence to the early search
    stage. GP posterior cross-covariance translates the regional preference
    into an adjustment of the mean used to compute expected improvement (EI).
    Standard EI is used when no valid regional preference is available.

    \item We evaluate LLMBO-MO on the Chen2020 and Ecker2015 battery
    parameterizations under limited simulation budgets. LLMBO-MO achieves a
    higher final mean hypervolume than ParEGO on both parameterizations at the
    tested budget. Initialization and ablation studies show that
    battery-specific initialization accounts for most of the improvement in
    the short-budget experiments.

\end{itemize}

The remainder of this paper is organized as follows.
Sections~\ref{sec:problem} and~\ref{sec:model_section} present the charging
problem and battery model, respectively. Section~\ref{sec:method} describes
LLMBO-MO. Section~\ref{sec:experiments} reports the experiments, and
Section~\ref{sec:conclusion} concludes the paper.
